\section{Build system and environment}

The build environment did not match the specification's assumptions on almost
every axis, and one mismatch was safety critical. The machine runs its Linux
toolchain under WSL2 with a deliberate resource cap of twelve gigabytes of memory
and eight processors, set after an earlier run without a cap exhausted host memory
and destabilized the machine. The specification suggested a larger budget and
twenty parallel compile jobs. Following that would have reintroduced the exact
condition that had caused the crash, so instead the one time LLVM and MLIR build
was configured with six compile jobs and a single link job, which is the guidance
in the resource cap file itself. The build completed in about twenty eight minutes
and stayed near three gigabytes of resident memory.

A second environment problem cost real time. The project folder path contained
spaces, and the LLVM lit and FileCheck test tools do not support spaces in paths;
they resolve symlinks, so a space free symlink over the spaced directory did not
help either. The fix was to host the repository in the WSL native filesystem,
which is also much faster to build in than the mounted Windows filesystem.

Finally, the entry points \texttt{registerAllDialects} and
\texttt{registerAllPasses} now live in their own aggregator libraries in this LLVM
release and are not pulled in by the dialect library properties, so the first link
of the optimizer driver failed with undefined references until those libraries
were listed explicitly.
